% Formalizing a Charting DSL: ANTLR4 Grammar Engineering and ASDL Intermediate
% Representations for Pine Script
% arXiv cs.PL language implementation paper (PYNE series, paper04-grammar)
\documentclass[11pt,a4paper]{article}

\usepackage[margin=1in]{geometry}
\usepackage[T1]{fontenc}
\usepackage[utf8]{inputenc}
\usepackage{lmodern}

% Shared macros for PYNE arXiv paper series
% Include with: % Shared macros for PYNE arXiv paper series
% Include with: % Shared macros for PYNE arXiv paper series
% Include with: \input{../common/macros.tex}

\usepackage{amsmath,amssymb,amsfonts}
\usepackage{amsthm}
\usepackage{booktabs}
\usepackage{graphicx}
\usepackage{xcolor}
\usepackage{hyperref}
\usepackage{url}
\usepackage{algorithm}
\usepackage{algpseudocode}
\usepackage{listings}
\usepackage{tikz}
\usepackage{pgfplots}
\pgfplotsset{compat=1.17}
\usetikzlibrary{arrows.meta,positioning,fit,calc,shapes.geometric,shapes.multipart,decorations.pathreplacing,backgrounds,chains,matrix,patterns}
\usepackage{subcaption}
\usepackage{multirow}
\usepackage{tabularx}
\usepackage{enumitem}
\usepackage{microtype}
\usepackage{balance}
\usepackage{csquotes}
\usepackage{natbib}

% Colors
\definecolor{pyneorange}{HTML}{F97316}
\definecolor{pyneblue}{HTML}{2563EB}
\definecolor{pynegray}{HTML}{6B7280}
\definecolor{pynegreen}{HTML}{16A34A}
\definecolor{codebg}{HTML}{F8FAFC}
\definecolor{codekw}{HTML}{7C3AED}
\definecolor{codestr}{HTML}{B45309}
\definecolor{codecomment}{HTML}{64748B}

\hypersetup{
  colorlinks=true,
  linkcolor=pyneblue,
  citecolor=pynegreen,
  urlcolor=pyneorange,
  pdftitle={},
  pdfauthor={HOOX / PYNE Project}
}

% Code listings
\lstdefinestyle{pine}{
  basicstyle=\ttfamily\footnotesize,
  backgroundcolor=\color{codebg},
  keywordstyle=\color{codekw}\bfseries,
  stringstyle=\color{codestr},
  commentstyle=\color{codecomment}\itshape,
  breaklines=true,
  frame=single,
  rulecolor=\color{black!20},
  numbers=left,
  numberstyle=\tiny\color{pynegray},
  numbersep=6pt,
  xleftmargin=1.5em,
  framexleftmargin=1.2em,
  showstringspaces=false,
  tabsize=2,
  morekeywords={indicator,strategy,library,plot,plotshape,fill,hline,var,varip,if,else,for,while,switch,true,false,na,and,or,not,import,export,type,method,enum},
  morecomment=[l]{//},
  morestring=[b]",
}
\lstdefinestyle{python}{
  language=Python,
  basicstyle=\ttfamily\footnotesize,
  backgroundcolor=\color{codebg},
  keywordstyle=\color{codekw}\bfseries,
  stringstyle=\color{codestr},
  commentstyle=\color{codecomment}\itshape,
  breaklines=true,
  frame=single,
  rulecolor=\color{black!20},
  numbers=left,
  numberstyle=\tiny\color{pynegray},
  numbersep=6pt,
  xleftmargin=1.5em,
  framexleftmargin=1.2em,
  showstringspaces=false,
  tabsize=4,
}
\lstset{style=python}

% Theorem environments
\theoremstyle{plain}
\newtheorem{theorem}{Theorem}[section]
\newtheorem{lemma}[theorem]{Lemma}
\newtheorem{proposition}[theorem]{Proposition}
\newtheorem{corollary}[theorem]{Corollary}
\theoremstyle{definition}
\newtheorem{definition}[theorem]{Definition}
\newtheorem{example}[theorem]{Example}
\newtheorem{remark}[theorem]{Remark}

% Shortcuts
\newcommand{\pyne}{\textsc{Pyne}}
\newcommand{\pine}{Pine Script\texttrademark}
\newcommand{\tv}{TradingView\textregistered}
\newcommand{\asdl}{\textsc{Asdl}}
\newcommand{\antlr}{\textsc{Antlr}4}
\newcommand{\numba}{\textsc{Numba}}
\newcommand{\lsp}{\textsc{Lsp}}
\newcommand{\ohlcv}{\textsc{Ohlcv}}
\newcommand{\udt}{\textsc{Udt}}
\newcommand{\jit}{\textsc{Jit}}
\newcommand{\api}{\textsc{Api}}
\newcommand{\cli}{\textsc{Cli}}
% AST acronym (text mode). Keep amsmath's math asterisk \ast intact.
\DeclareRobustCommand{\Ast}{\textsc{Ast}}
% Papers in this series write \ast{} for the AST acronym in prose.
% Redefine robustly; use \mathast for a binary asterisk in formulas.
\let\mathast\ast
\DeclareRobustCommand{\ast}{\textsc{Ast}}
\newcommand{\ir}{\textsc{Ir}}
\newcommand{\ta}{\texttt{ta}}
\newcommand{\code}[1]{\texttt{#1}}
\newcommand{\file}[1]{\texttt{#1}}
\newcommand{\eg}{e.g.}
\newcommand{\ie}{i.e.}
\newcommand{\etal}{\textit{et~al.}}
\newcommand{\resp}{respectively}

% Math helpers
\newcommand{\R}{\mathbb{R}}
\newcommand{\N}{\mathbb{N}}
\newcommand{\Z}{\mathbb{Z}}
% Text-mode NA token for prose; use $\namath$ inside formulas if needed.
\DeclareRobustCommand{\na}{\textsf{na}}
\newcommand{\namath}{\mathsf{na}}
\newcommand{\series}[1]{\mathbf{#1}}
\newcommand{\baridx}{i}
\newcommand{\bars}{n}
\newcommand{\period}{p}
\newcommand{\abserr}{\varepsilon_{\mathrm{abs}}}
\newcommand{\relerr}{\varepsilon_{\mathrm{rel}}}
\newcommand{\maxerr}{\varepsilon_{\mathrm{max}}}
\newcommand{\meanerr}{\varepsilon_{\mathrm{mean}}}

% Pipeline notation — use inside math: $\pipeline$
\newcommand{\pipeline}{\mathrm{src}\!\to\!\mathrm{lex}\!\to\!\mathrm{parse}\!\to\!\mathrm{AST}\!\to\!\mathrm{eval}}

% Disclaimer footnote helper
\newcommand{\trademarknote}{%
  \pine{} and \tv{} are trademarks of TradingView, Inc.
  \pyne{} is an independent, unofficial implementation and is not affiliated with,
  authorized by, sponsored by, or endorsed by TradingView, Inc.
}

% Figure path helper (papers live one level under .papers/)
\graphicspath{{figures/}{../common/figures/}}


\usepackage{amsmath,amssymb,amsfonts}
\usepackage{amsthm}
\usepackage{booktabs}
\usepackage{graphicx}
\usepackage{xcolor}
\usepackage{hyperref}
\usepackage{url}
\usepackage{algorithm}
\usepackage{algpseudocode}
\usepackage{listings}
\usepackage{tikz}
\usepackage{pgfplots}
\pgfplotsset{compat=1.17}
\usetikzlibrary{arrows.meta,positioning,fit,calc,shapes.geometric,shapes.multipart,decorations.pathreplacing,backgrounds,chains,matrix,patterns}
\usepackage{subcaption}
\usepackage{multirow}
\usepackage{tabularx}
\usepackage{enumitem}
\usepackage{microtype}
\usepackage{balance}
\usepackage{csquotes}
\usepackage{natbib}

% Colors
\definecolor{pyneorange}{HTML}{F97316}
\definecolor{pyneblue}{HTML}{2563EB}
\definecolor{pynegray}{HTML}{6B7280}
\definecolor{pynegreen}{HTML}{16A34A}
\definecolor{codebg}{HTML}{F8FAFC}
\definecolor{codekw}{HTML}{7C3AED}
\definecolor{codestr}{HTML}{B45309}
\definecolor{codecomment}{HTML}{64748B}

\hypersetup{
  colorlinks=true,
  linkcolor=pyneblue,
  citecolor=pynegreen,
  urlcolor=pyneorange,
  pdftitle={},
  pdfauthor={HOOX / PYNE Project}
}

% Code listings
\lstdefinestyle{pine}{
  basicstyle=\ttfamily\footnotesize,
  backgroundcolor=\color{codebg},
  keywordstyle=\color{codekw}\bfseries,
  stringstyle=\color{codestr},
  commentstyle=\color{codecomment}\itshape,
  breaklines=true,
  frame=single,
  rulecolor=\color{black!20},
  numbers=left,
  numberstyle=\tiny\color{pynegray},
  numbersep=6pt,
  xleftmargin=1.5em,
  framexleftmargin=1.2em,
  showstringspaces=false,
  tabsize=2,
  morekeywords={indicator,strategy,library,plot,plotshape,fill,hline,var,varip,if,else,for,while,switch,true,false,na,and,or,not,import,export,type,method,enum},
  morecomment=[l]{//},
  morestring=[b]",
}
\lstdefinestyle{python}{
  language=Python,
  basicstyle=\ttfamily\footnotesize,
  backgroundcolor=\color{codebg},
  keywordstyle=\color{codekw}\bfseries,
  stringstyle=\color{codestr},
  commentstyle=\color{codecomment}\itshape,
  breaklines=true,
  frame=single,
  rulecolor=\color{black!20},
  numbers=left,
  numberstyle=\tiny\color{pynegray},
  numbersep=6pt,
  xleftmargin=1.5em,
  framexleftmargin=1.2em,
  showstringspaces=false,
  tabsize=4,
}
\lstset{style=python}

% Theorem environments
\theoremstyle{plain}
\newtheorem{theorem}{Theorem}[section]
\newtheorem{lemma}[theorem]{Lemma}
\newtheorem{proposition}[theorem]{Proposition}
\newtheorem{corollary}[theorem]{Corollary}
\theoremstyle{definition}
\newtheorem{definition}[theorem]{Definition}
\newtheorem{example}[theorem]{Example}
\newtheorem{remark}[theorem]{Remark}

% Shortcuts
\newcommand{\pyne}{\textsc{Pyne}}
\newcommand{\pine}{Pine Script\texttrademark}
\newcommand{\tv}{TradingView\textregistered}
\newcommand{\asdl}{\textsc{Asdl}}
\newcommand{\antlr}{\textsc{Antlr}4}
\newcommand{\numba}{\textsc{Numba}}
\newcommand{\lsp}{\textsc{Lsp}}
\newcommand{\ohlcv}{\textsc{Ohlcv}}
\newcommand{\udt}{\textsc{Udt}}
\newcommand{\jit}{\textsc{Jit}}
\newcommand{\api}{\textsc{Api}}
\newcommand{\cli}{\textsc{Cli}}
% AST acronym (text mode). Keep amsmath's math asterisk \ast intact.
\DeclareRobustCommand{\Ast}{\textsc{Ast}}
% Papers in this series write \ast{} for the AST acronym in prose.
% Redefine robustly; use \mathast for a binary asterisk in formulas.
\let\mathast\ast
\DeclareRobustCommand{\ast}{\textsc{Ast}}
\newcommand{\ir}{\textsc{Ir}}
\newcommand{\ta}{\texttt{ta}}
\newcommand{\code}[1]{\texttt{#1}}
\newcommand{\file}[1]{\texttt{#1}}
\newcommand{\eg}{e.g.}
\newcommand{\ie}{i.e.}
\newcommand{\etal}{\textit{et~al.}}
\newcommand{\resp}{respectively}

% Math helpers
\newcommand{\R}{\mathbb{R}}
\newcommand{\N}{\mathbb{N}}
\newcommand{\Z}{\mathbb{Z}}
% Text-mode NA token for prose; use $\namath$ inside formulas if needed.
\DeclareRobustCommand{\na}{\textsf{na}}
\newcommand{\namath}{\mathsf{na}}
\newcommand{\series}[1]{\mathbf{#1}}
\newcommand{\baridx}{i}
\newcommand{\bars}{n}
\newcommand{\period}{p}
\newcommand{\abserr}{\varepsilon_{\mathrm{abs}}}
\newcommand{\relerr}{\varepsilon_{\mathrm{rel}}}
\newcommand{\maxerr}{\varepsilon_{\mathrm{max}}}
\newcommand{\meanerr}{\varepsilon_{\mathrm{mean}}}

% Pipeline notation — use inside math: $\pipeline$
\newcommand{\pipeline}{\mathrm{src}\!\to\!\mathrm{lex}\!\to\!\mathrm{parse}\!\to\!\mathrm{AST}\!\to\!\mathrm{eval}}

% Disclaimer footnote helper
\newcommand{\trademarknote}{%
  \pine{} and \tv{} are trademarks of TradingView, Inc.
  \pyne{} is an independent, unofficial implementation and is not affiliated with,
  authorized by, sponsored by, or endorsed by TradingView, Inc.
}

% Figure path helper (papers live one level under .papers/)
\graphicspath{{figures/}{../common/figures/}}


\usepackage{amsmath,amssymb,amsfonts}
\usepackage{amsthm}
\usepackage{booktabs}
\usepackage{graphicx}
\usepackage{xcolor}
\usepackage{hyperref}
\usepackage{url}
\usepackage{algorithm}
\usepackage{algpseudocode}
\usepackage{listings}
\usepackage{tikz}
\usepackage{pgfplots}
\pgfplotsset{compat=1.17}
\usetikzlibrary{arrows.meta,positioning,fit,calc,shapes.geometric,shapes.multipart,decorations.pathreplacing,backgrounds,chains,matrix,patterns}
\usepackage{subcaption}
\usepackage{multirow}
\usepackage{tabularx}
\usepackage{enumitem}
\usepackage{microtype}
\usepackage{balance}
\usepackage{csquotes}
\usepackage{natbib}

% Colors
\definecolor{pyneorange}{HTML}{F97316}
\definecolor{pyneblue}{HTML}{2563EB}
\definecolor{pynegray}{HTML}{6B7280}
\definecolor{pynegreen}{HTML}{16A34A}
\definecolor{codebg}{HTML}{F8FAFC}
\definecolor{codekw}{HTML}{7C3AED}
\definecolor{codestr}{HTML}{B45309}
\definecolor{codecomment}{HTML}{64748B}

\hypersetup{
  colorlinks=true,
  linkcolor=pyneblue,
  citecolor=pynegreen,
  urlcolor=pyneorange,
  pdftitle={},
  pdfauthor={HOOX / PYNE Project}
}

% Code listings
\lstdefinestyle{pine}{
  basicstyle=\ttfamily\footnotesize,
  backgroundcolor=\color{codebg},
  keywordstyle=\color{codekw}\bfseries,
  stringstyle=\color{codestr},
  commentstyle=\color{codecomment}\itshape,
  breaklines=true,
  frame=single,
  rulecolor=\color{black!20},
  numbers=left,
  numberstyle=\tiny\color{pynegray},
  numbersep=6pt,
  xleftmargin=1.5em,
  framexleftmargin=1.2em,
  showstringspaces=false,
  tabsize=2,
  morekeywords={indicator,strategy,library,plot,plotshape,fill,hline,var,varip,if,else,for,while,switch,true,false,na,and,or,not,import,export,type,method,enum},
  morecomment=[l]{//},
  morestring=[b]",
}
\lstdefinestyle{python}{
  language=Python,
  basicstyle=\ttfamily\footnotesize,
  backgroundcolor=\color{codebg},
  keywordstyle=\color{codekw}\bfseries,
  stringstyle=\color{codestr},
  commentstyle=\color{codecomment}\itshape,
  breaklines=true,
  frame=single,
  rulecolor=\color{black!20},
  numbers=left,
  numberstyle=\tiny\color{pynegray},
  numbersep=6pt,
  xleftmargin=1.5em,
  framexleftmargin=1.2em,
  showstringspaces=false,
  tabsize=4,
}
\lstset{style=python}

% Theorem environments
\theoremstyle{plain}
\newtheorem{theorem}{Theorem}[section]
\newtheorem{lemma}[theorem]{Lemma}
\newtheorem{proposition}[theorem]{Proposition}
\newtheorem{corollary}[theorem]{Corollary}
\theoremstyle{definition}
\newtheorem{definition}[theorem]{Definition}
\newtheorem{example}[theorem]{Example}
\newtheorem{remark}[theorem]{Remark}

% Shortcuts
\newcommand{\pyne}{\textsc{Pyne}}
\newcommand{\pine}{Pine Script\texttrademark}
\newcommand{\tv}{TradingView\textregistered}
\newcommand{\asdl}{\textsc{Asdl}}
\newcommand{\antlr}{\textsc{Antlr}4}
\newcommand{\numba}{\textsc{Numba}}
\newcommand{\lsp}{\textsc{Lsp}}
\newcommand{\ohlcv}{\textsc{Ohlcv}}
\newcommand{\udt}{\textsc{Udt}}
\newcommand{\jit}{\textsc{Jit}}
\newcommand{\api}{\textsc{Api}}
\newcommand{\cli}{\textsc{Cli}}
% AST acronym (text mode). Keep amsmath's math asterisk \ast intact.
\DeclareRobustCommand{\Ast}{\textsc{Ast}}
% Papers in this series write \ast{} for the AST acronym in prose.
% Redefine robustly; use \mathast for a binary asterisk in formulas.
\let\mathast\ast
\DeclareRobustCommand{\ast}{\textsc{Ast}}
\newcommand{\ir}{\textsc{Ir}}
\newcommand{\ta}{\texttt{ta}}
\newcommand{\code}[1]{\texttt{#1}}
\newcommand{\file}[1]{\texttt{#1}}
\newcommand{\eg}{e.g.}
\newcommand{\ie}{i.e.}
\newcommand{\etal}{\textit{et~al.}}
\newcommand{\resp}{respectively}

% Math helpers
\newcommand{\R}{\mathbb{R}}
\newcommand{\N}{\mathbb{N}}
\newcommand{\Z}{\mathbb{Z}}
% Text-mode NA token for prose; use $\namath$ inside formulas if needed.
\DeclareRobustCommand{\na}{\textsf{na}}
\newcommand{\namath}{\mathsf{na}}
\newcommand{\series}[1]{\mathbf{#1}}
\newcommand{\baridx}{i}
\newcommand{\bars}{n}
\newcommand{\period}{p}
\newcommand{\abserr}{\varepsilon_{\mathrm{abs}}}
\newcommand{\relerr}{\varepsilon_{\mathrm{rel}}}
\newcommand{\maxerr}{\varepsilon_{\mathrm{max}}}
\newcommand{\meanerr}{\varepsilon_{\mathrm{mean}}}

% Pipeline notation — use inside math: $\pipeline$
\newcommand{\pipeline}{\mathrm{src}\!\to\!\mathrm{lex}\!\to\!\mathrm{parse}\!\to\!\mathrm{AST}\!\to\!\mathrm{eval}}

% Disclaimer footnote helper
\newcommand{\trademarknote}{%
  \pine{} and \tv{} are trademarks of TradingView, Inc.
  \pyne{} is an independent, unofficial implementation and is not affiliated with,
  authorized by, sponsored by, or endorsed by TradingView, Inc.
}

% Figure path helper (papers live one level under .papers/)
\graphicspath{{figures/}{../common/figures/}}


\hypersetup{
  pdftitle={Formalizing a Charting DSL: ANTLR4 Grammar Engineering and ASDL Intermediate Representations for Pine Script},
  pdfauthor={HOOX Research / PYNE Contributors},
  pdfkeywords={grammar engineering, ANTLR, ASDL, abstract syntax trees, parse-unparse round-trip, domain-specific languages, Pine Script}
}

\title{%
  Formalizing a Charting DSL:\\
  \large ANTLR4 Grammar Engineering and ASDL Intermediate Representations\\
  for \pine{}
}
\author{%
  HOOX Research / \pyne{} Contributors\\
  \texttt{jango\_blockchained}\\
  Independent Open Source Research\\
  \texttt{contact@hoox.sh}
}
\date{}

\begin{document}
\maketitle

\begin{abstract}
Domain-specific languages (DSLs) for financial charting combine expression-oriented
semantics with significant whitespace, versioned surface syntax, and a dense
builtin surface inventory. Building reliable tooling---parsers, formatters, linters,
language servers, and interpreters---requires a front-end that is both faithful to
community scripts and amenable to static analysis.
We present the language front-end of \pyne{} (package \texttt{hoox-pyne}), an
independent open toolchain for \pine{}: a hand-engineered \antlr{} grammar with
Python-side indent tracking, a Zephyr-\asdl{} abstract syntax tree (AST) schema,
a visitor-based parse-tree builder, a precedence-aware unparser, a structural
linter, and a qualifier-aware type model.
The pipeline is
\[
  \text{source}
  \to \text{FileStream/InputStream}
  \to \text{Lexer}
  \to \text{CommonTokenStream}
  \to \text{Parser}
  \to \text{parse tree}
  \to \text{AST builder}
  \to \asdl{}\ \ast
  \to \text{annotations}
  \to \text{unparse}.
\]
Hand-edited artifacts are limited to \file{antlr4/resource/*.g4} and
\file{asdl/resource/Pinescript.asdl}; generated lexer/parser and node classes are
committed so end users need no Java toolchain.
We detail SLL-first / LL-fallback prediction (\(\approx 5.4\times\) on complex
scripts), annotation attachment for \code{//@version} and related markers,
parse-cache design (SHA-256 LRU with call-site scrubbing; multi-parse
\(\approx 1000\times\) on hits), unparser thread-local reuse
(\(\approx 1.95\times\) on a 99-script corpus), corpus sanitization for scraped
sources, and evaluation against \(138+\) real scripts with \(1100+\) automated
tests in the broader project.
The grammar is approximate and unofficial: \pyne{} is not affiliated with
TradingView, Inc.
\end{abstract}

\paragraph{Keywords.}
grammar engineering, ANTLR, ASDL, abstract syntax trees, parse--unparse
round-trip, domain-specific languages, \pine{}

\paragraph{Disclaimer.}
\trademarknote{}

%==============================================================================
\section{Introduction}
\label{sec:intro}
%==============================================================================

Charting and algorithmic trading communities rely on specialized languages that
embed market data series, drawing primitives, and strategy execution into a
compact surface syntax~\citep{fowler2010dsl,hudak1996dsl,tradingview-pine}.
\pine{} is one such language: scripts are short relative to general-purpose
codebases yet dense in temporal operators, qualifiers (\code{const},
\code{simple}, \code{series}, \code{input}), user-defined types (UDTs), methods,
and (from v6) enums and multiline strings.
Tooling that merely tokenizes for editor highlighting is insufficient for
format-on-save, structural search, type-directed completion, or faithful
interpretation~\citep{lsp2023,nystrom2021crafting}.

\pyne{}~\citep{pynescript2026} is an independent, AGPL-licensed implementation of a
\pine{} toolchain in Python. This paper isolates the \emph{language front-end}:
everything from source text to a stable \asdl{} tree and back, including the
grammar engineering practices that make the rest of the stack possible.
Companion papers in the series address overall architecture, compilation to
\numba{}, series runtime semantics, strategy parity, and LSP integration; here
we focus on:

\begin{enumerate}[leftmargin=*]
  \item \textbf{Grammar design} for a whitespace-significant, expression-oriented
    charting DSL, including indent tokens, string forms, colors, and v6 surface
    extensions.
  \item \textbf{\asdl{} intermediate representation} as a single source of truth
    for node shapes, shared by the builder, unparser, linter, evaluator, and
    compiler.
  \item \textbf{Construction and unparsing} with annotation attachment and
    precedence-aware pretty-printing.
  \item \textbf{Static analysis hooks}: a rule-based linter and a type model for
    qualifiers, collections, and UDTs.
  \item \textbf{Corpus engineering}: sanitization of scraped community sources
    and a regeneration workflow that keeps generated artifacts reproducible.
  \item \textbf{Evaluation}: parse success, round-trip stability, and
    performance tables for SLL/LL, unparse, and cache warm paths.
\end{enumerate}

\paragraph{Contributions.}
\begin{itemize}[leftmargin=*]
  \item A documented \antlr{} grammar and lexer base for \pine{} v5/v6 surface
    features, with generated Python targets committed for zero-Java installs.
  \item An \asdl{} schema and visitor infrastructure analogous to CPython's
    \code{ast} module~\citep{asdl-python,wang1997asdl}.
  \item A lossless-enough round-trip for structure and \code{//@}-annotations,
    powering CLI format and LSP document formatting.
  \item Measured front-end optimizations (SLL-first parse, annotation skip,
    thread-local unparser, SHA-256 parse cache) with correctness checks via
    AST dump identity and re-parse stability.
\end{itemize}

%==============================================================================
\section{Background}
\label{sec:background}
%==============================================================================

\subsection{Pine surface evolution (v1--v6)}
\label{sec:pine-evolution}

\pine{} evolved from a narrow indicator dialect into a multiparadigm charting
language~\citep{tradingview-pine}. Early versions emphasized single-line
expressions and a small set of drawing builtins. Later major versions introduced:

\begin{description}[style=nextline,leftmargin=1.2em]
  \item[v3--v4]
    Multi-line scripts, \code{var} persistence, expanded \code{security}
    patterns, and broader strategy APIs.
  \item[v5]
    Namespaced builtins (\code{ta.*}, \code{math.*}, \code{strategy.*}, \ldots),
    libraries (\code{import} / \code{export}), methods, UDTs (\code{type}),
    \code{switch}, and richer type annotations on parameters.
  \item[v6]
    \code{enum} declarations; triple-quoted multiline strings
    (\code{"""\ldots"""} / \code{'''\ldots'''}); stricter boolean typing;
    integer division producing floats in more contexts; collection helpers such
    as \code{sort\_field} on \code{array}/\code{matrix} sort; and continued
    library-surface growth.
\end{description}

Community scripts frequently mix versions, rely on soft conventions (camelCase
identifiers, trailing commas), and embed annotation comments
(\code{//@version=6}, \code{//@description}, \code{//@function}) that are not
ordinary statements yet must survive tooling passes.

\subsection{Challenges of a whitespace-significant, expression-oriented DSL}
\label{sec:challenges}

From a language-implementation perspective~\citep{aho2006dragon,cooper2012engineering,nystrom2021crafting},
\pine{} combines several properties that stress classical parser generators:

\begin{enumerate}[leftmargin=*]
  \item \textbf{Indentation as structure.}
    Blocks after \code{if}/\code{for}/\code{while}/\code{switch}/\code{type}/
    \code{enum}/\code{=>} use Python-like INDENT/DEDENT tokens. Indent width is
    conventionally four spaces; line wrapping may use non-multiple-of-four
    indents that must \emph{not} introduce structure.
  \item \textbf{Structures as expressions.}
    \code{if}, \code{for}, \code{switch}, and related forms appear both as
    statements and as right-hand sides of assignments (``structure
    expressions''), so the grammar cannot segregate expression and statement
    nonterminals as cleanly as C or Java.
  \item \textbf{Soft keywords and contextual tokens.}
    Names such as type constructors interact with primary-expression suffixes
    (attributes, calls, subscripts, template specs like \code{array.new<float>}).
  \item \textbf{Literal diversity.}
    Numbers (decimal/binary/octal/hex, floats, optional underscores), colors
    (\code{\#RRGGBB}, \code{\#RRGGBBAA}), and strings (escaped single-line and
    v6 triple-quoted multiline) coexist with comment channels and paste noise.
  \item \textbf{Versioned semantics outside pure syntax.}
    Boolean strictness and division typing are semantic; the front-end must
    still accept the surface forms so later stages can apply version policy.
  \item \textbf{Scraped corpora.}
    Real evaluation sets often arrive wrapped in markdown fences, UI chrome
    (``Expand (\(N\) lines)''), prose, and HTML---not clean \code{.pine} files.
\end{enumerate}

These challenges motivate a hybrid design: ANTLR for adaptive \(\mathrm{LL}(*)\)
parsing~\citep{parr2013antlr,parr2011allstar}, a hand-written lexer base for
indent and newline policy, and an \asdl{} IR that decouples consumers from
parse-tree volatility~\citep{wang1997asdl,asdl-python}.

%==============================================================================
\section{Grammar Design}
\label{sec:grammar}
%==============================================================================

\subsection{Artifact layout and generation}
\label{sec:artifact-layout}

The front-end distinguishes \emph{hand-edited resources} from
\emph{generated, committed} products:

\begin{center}
\small
\begin{tabular}{@{}llp{0.48\textwidth}@{}}
\toprule
Kind & Path (under \file{src/pynescript/ast/grammar/}) & Role \\
\midrule
Hand & \file{antlr4/resource/PinescriptLexer.g4} & Tokens, fragments, channels \\
Hand & \file{antlr4/resource/PinescriptParser.g4} & Syntactic productions \\
Hand & \file{antlr4/resource/PinescriptLexerBase.py} & Indent / newline engine \\
Hand & \file{antlr4/resource/PinescriptParserBase.py} & Parser super-class hooks \\
Hand & \file{asdl/resource/Pinescript.asdl} & Node schema \\
Generated & \file{antlr4/generated/*} & Python lexer/parser/visitor \\
Generated & \file{asdl/generated/*} & Node classes \\
\bottomrule
\end{tabular}
\end{center}

Regeneration uses project tooling (e.g.\ \code{hatch run lint:gen-parser});
committing generated Python removes the Java/ANTLR runtime dependency for
ordinary installs~\citep{parr2013antlr,antlr4}.

Figure~\ref{fig:pipeline} summarizes the end-to-end front-end pipeline.

\begin{figure}[t]
\centering
\begin{tikzpicture}[
  font=\small,
  box/.style={draw, rounded corners=2pt, align=center, inner sep=5pt,
    minimum height=0.85cm, fill=codebg},
  arr/.style={-{Latex}, thick, pyneblue}
]
\node[box] (src) {Pine source};
\node[box, right=0.45cm of src] (fs) {InputStream\\/FileStream};
\node[box, right=0.45cm of fs] (lex) {PinescriptLexer\\(+ LexerBase)};
\node[box, right=0.45cm of lex] (cts) {CommonToken\\Stream};
\node[box, below=0.7cm of cts] (par) {PinescriptParser\\SLL \(\to\) LL};
\node[box, left=0.45cm of par] (pt) {Parse tree};
\node[box, left=0.45cm of pt] (bld) {ASTBuilder\\visit\_*};
\node[box, left=0.45cm of bld] (astn) {\asdl{} AST\\+ annotations};
\node[box, below=0.7cm of astn] (unp) {Unparser};
\node[box, right=0.45cm of unp] (out) {Pine source};
\draw[arr] (src) -- (fs);
\draw[arr] (fs) -- (lex);
\draw[arr] (lex) -- (cts);
\draw[arr] (cts) -- (par);
\draw[arr] (par) -- (pt);
\draw[arr] (pt) -- (bld);
\draw[arr] (bld) -- (astn);
\draw[arr] (astn) -- (unp);
\draw[arr] (unp) -- (out);
\draw[arr, dashed, pyneorange] (out.north) to[bend left=20] node[above, font=\scriptsize, text=pynegray]{round-trip} (src.south);
\end{tikzpicture}
\caption{Front-end pipeline: lex \(\to\) parse \(\to\) build \(\to\) annotate
\(\to\) unparse. Solid arrows form the primary data path; the dashed arrow
denotes structural round-trip used by formatters and tests.}
\label{fig:pipeline}
\end{figure}

\subsection{Lexical structure}
\label{sec:lexical}

\subsubsection{Keywords, operators, and channels}
The lexer grammar declares keywords (\code{and}, \code{or}, \code{var},
\code{varip}, \code{type}, \code{method}, \code{enum}, \ldots), multi-character
operators before single-character ones (so longest-match prefers \code{<=} over
\code{<}), bitwise operators (\code{\&}, \code{|}, \code{\^{}}, \code{<<},
\code{>>}, \code{\~}), and assignment forms (\code{=}, \code{:=},
\code{+=}, \ldots).
Comments travel on \code{COMMENT\_CHANNEL}; whitespace is \code{HIDDEN}.
Additional robustness rules accept \code{\#}-style line comments when not
part of a color literal, markdown backticks as comment noise, and selected
Unicode punctuation as hidden ``paste noise.''

\subsubsection{Indentation tracking}
\code{PinescriptLexerBase} (Python super-class of the generated lexer)
implements the indent engine:

\begin{itemize}[leftmargin=*]
  \item Stack of indent lengths; emit synthetic \code{INDENT}/\code{DEDENT}
    tokens when structure changes.
  \item Ignore newlines inside open \code{(}/\code{[}, after operators (line
    continuation), and for wraps whose indent is not a multiple of four.
  \item Collapse consecutive newlines; ensure a trailing newline at EOF.
  \item Special handling of multiline string literals so wrap indentation does
    not become structural.
\end{itemize}

This design mirrors classical off-side-rule lexers~\citep{aho2006dragon} while
remaining local to the ANTLR token stream.

\subsubsection{Strings, colors, numbers, identifiers}
Listing~\ref{lst:lexer-strings} excerpts the v6 triple-quote rules and color
fragments. Numbers support base prefixes and digit separators; identifiers use
Unicode letter classes (\code{\textbackslash p\{L\}}) so non-ASCII names in
community scripts parse cleanly.

\begin{lstlisting}[style=pine,caption={Lexer excerpts: triple-quoted strings and color literals (simplified from \texttt{PinescriptLexer.g4}).},label={lst:lexer-strings},language={}]
// v6 multiline triple-quoted strings
TRIPLE_DQ_STRING : '"""' ( ~'"' | '"' ~'"' | '""' ~'"' )* '"""' -> type(STRING) ;
TRIPLE_SQ_STRING : '\'\'\'' ( ~'\'' | '\'' ~'\'' | '\'\'' ~'\'' )* '\'\'\'' -> type(STRING) ;

// #RRGGBB or #RRGGBBAA (HASH_COMMENT must not steal these)
fragment COLOR_LITERAL_RGBA : '#' HEX_DIGIT{8} ;
fragment COLOR_LITERAL_RGB  : '#' HEX_DIGIT{6} ;
\end{lstlisting}

\subsection{Syntactic productions}
\label{sec:syntax}

\subsubsection{Scripts and statement forms}
Start rules are \code{start\_script} (full program, mode \code{exec}) and
\code{start\_expression} (mode \code{eval}). A script is a sequence of
statements: compound forms (assignments with structure RHSs, type/enum/method/
function declarations, control structures) and simple statements (assignments,
expressions, imports, \code{break}/\code{continue}), optionally packed on one
line with commas. A distinctive production,
\code{trailing\_structure\_statements}, accepts real-world lines such as
\begin{lstlisting}[style=pine,numbers=none]
Ex = 0.0, Ey = 0.0, for i = 0 to n
    ...
\end{lstlisting}

\subsubsection{Declarations: functions, methods, types, enums}
Function and method declarations use \code{=>} and a local block (indented or
inline). Optional leading type specifications encode return types
(\code{int ilog2(int n) => \ldots}). Type declarations introduce UDT fields under
INDENT/DEDENT; enum declarations introduce member definitions with optional
values. \code{export} prefixes support library authors.

\subsubsection{Expressions}
Expressions are layered by precedence (Figure~\ref{fig:syntax-overview}):
conditional (ternary), disjunction, conjunction, bitwise OR/XOR/AND, equality,
inequality, shifts, additive, multiplicative, unary, and primary
(attribute/call/subscript/template). Logical \code{and} binds less tightly than
bitwise OR, matching Pine's C-like bitwise layering. Primary expressions use
labeled alternatives for attribute, call, and subscript forms.

\begin{figure}[t]
\centering
\begin{tikzpicture}[
  font=\scriptsize,
  level 1/.style={sibling distance=3.6cm, level distance=1.1cm},
  level 2/.style={sibling distance=1.7cm, level distance=1.0cm},
  level 3/.style={sibling distance=1.1cm, level distance=0.9cm},
  every node/.style={draw, rounded corners=1pt, align=center, fill=codebg, inner sep=2pt}
]
\node {script}
  child { node {stmt}
    child { node {compound\\decl/ctrl} }
    child { node {simple\\assign/expr} }
    child { node {trailing\\structure} }
  }
  child { node {expr}
    child { node {ternary} }
    child { node {bool/bitwise} }
    child { node {arith/compare} }
    child { node {primary\\attr/call/[]} }
  }
  child { node {type forms}
    child { node {UDT} }
    child { node {enum} }
    child { node {method} }
  };
\end{tikzpicture}
\caption{Hierarchical overview of major syntactic families in
\file{PinescriptParser.g4}. Structures double as expressions on assignment RHSs.}
\label{fig:syntax-overview}
\end{figure}

Listing~\ref{lst:g4-fn} shows the function and enum productions.

\begin{lstlisting}[style=pine,caption={Parser productions for functions and enums (from \texttt{PinescriptParser.g4}).},label={lst:g4-fn},language={}]
function_declaration
  : EXPORT? type_specification? name LPAR parameter_list? RPAR RARROW local_block ;

enum_declaration
  : EXPORT? ENUM name NEWLINE INDENT enum_definitions DEDENT ;

enum_definition: name_store (EQUAL expression)? NEWLINE ;
\end{lstlisting}

\subsection{SLL / LL prediction strategy}
\label{sec:sll}

ANTLR's adaptive \(\mathrm{LL}(*)\) engine can fall into expensive full-context
prediction on ambiguous decision points~\citep{parr2011allstar,parr2013antlr}.
\pyne{} uses the standard two-stage strategy in \file{helper.py}:

\begin{enumerate}[leftmargin=*]
  \item Parse with \code{PredictionMode.SLL} and \code{BailErrorStrategy}.
  \item On \code{ParseCancellationException}: \code{token\_stream.seek(0)},
    \code{parser.reset()}, switch to \code{PredictionMode.LL} with
    \code{DefaultErrorStrategy}, and re-parse.
\end{enumerate}

On stratified mid-size scripts, SLL-first produced identical AST dumps to
pure LL with zero fallbacks in the measured sample, while reducing wall time by
approximately \(5.4\times\) (Section~\ref{sec:eval}). Strategy instances and a
shared builder are reused to avoid per-parse allocation.

\subsection{Error listeners}
\label{sec:errors}

\code{PinescriptErrorListener} is a singleton attached to both lexer and
parser after default listeners are removed. Syntax failures surface as
\code{pynescript.ast.error.SyntaxError} with optional
\code{SyntaxErrorDetails} (filename, line, column, source excerpt). Indentation
violations raise a dedicated \code{IndentationError}. The linter maps these
structured locations into diagnostic chips for editors (Section~\ref{sec:static}).

%==============================================================================
\section{ASDL Intermediate Representation}
\label{sec:asdl}
%==============================================================================

\subsection{Schema philosophy}
\label{sec:asdl-philosophy}

Zephyr ASDL~\citep{wang1997asdl} describes tree shapes as sum and product types.
Python's compiler has long used ASDL for the \code{ast} module~\citep{asdl-python};
\pyne{} follows the same pattern so that:

\begin{itemize}[leftmargin=*]
  \item Node classes are generated, not hand-synchronized.
  \item Location attributes (\code{lineno}, \code{col\_offset},
    \code{end\_lineno}, \code{end\_col\_offset}) are uniform.
  \item Visitors and transformers can be written once against a stable API.
\end{itemize}

The module root distinguishes full scripts from expression-mode trees:

\begin{lstlisting}[style=python,caption={ASDL module root and statement taxonomy (excerpt from \texttt{Pinescript.asdl}).},label={lst:asdl}]
module Pinescript
{
  mod = Script(stmt* body, string* annotations)
      | Expression(expr body)

  stmt = FunctionDef(identifier name, param* args, stmt* body,
                     int? method, int? export, string* annotations)
       | TypeDef(identifier name, stmt* body, int? export, string* annotations)
       | EnumDef(identifier name, stmt* body, int? export, string* annotations)
       | Assign(expr target, expr? value, expr? type, decl_mode? mode,
                int? export, string* annotations)
       | ReAssign(expr target, expr value)
       | AugAssign(expr target, operator op, expr value)
       | Import(identifier namespace, identifier name, int version,
                identifier? alias)
       | Expr(expr value) | Break | Continue
       attributes (int lineno, int col_offset,
                   int? end_lineno, int? end_col_offset)
  ...
}
\end{lstlisting}

\subsection{Node taxonomy}
\label{sec:taxonomy}

Figure~\ref{fig:asdl-tree} sketches the conceptual hierarchy.

\begin{figure}[t]
\centering
\begin{tikzpicture}[
  font=\small,
  edge from parent/.style={draw, -{Latex}, thick, pyneblue},
  level 1/.style={sibling distance=4.2cm, level distance=1.25cm},
  level 2/.style={sibling distance=1.55cm, level distance=1.15cm},
  every node/.style={draw, rounded corners=2pt, align=center, fill=white,
    inner sep=3pt, minimum height=0.55cm}
]
\node[fill=pyneorange!15] {AST}
  child { node[fill=pyneblue!12] {mod}
    child { node {Script} }
    child { node {Expression} }
  }
  child { node[fill=pynegreen!12] {stmt}
    child { node {Assign\\ReAssign} }
    child { node {FunctionDef\\Method} }
    child { node {TypeDef\\EnumDef} }
    child { node {Import\\Ctrl} }
  }
  child { node[fill=codebg] {expr}
    child { node {BinOp\\BoolOp} }
    child { node {Call\\Name} }
    child { node {If/For/\\Switch} }
    child { node {Qualify\\Specialize} }
  };
\end{tikzpicture}
\caption{\asdl{} node hierarchy (conceptual). Control structures are
expression nodes when used as values; declarations live under \texttt{stmt}.
Auxiliary products include \texttt{param}, \texttt{arg}, \texttt{case}, and
\texttt{Comment}.}
\label{fig:asdl-tree}
\end{figure}

\paragraph{Expressions.}
Algebraic nodes (\code{BoolOp}, \code{BinOp}, \code{UnaryOp}, \code{Compare},
\code{Conditional}) coexist with \code{Call}, \code{Attribute},
\code{Subscript}, \code{Name}, \code{Tuple}, and \code{Constant}. Control forms
\code{If}, \code{ForTo}, \code{ForIn}, \code{While}, and \code{Switch} are
\emph{expressions} in the schema, reflecting Pine's expression-oriented
structures. Type constructors use \code{Qualify} (const/input/simple/series)
and \code{Specialize} (generic arguments).

\paragraph{Enums of operators and modes.}
\code{decl\_mode} \(\in \{\code{Var},\code{VarIp}\}\);
\code{type\_qual} \(\in \{\code{Const},\code{Input},\code{Simple},\code{Series}\}\);
\code{expr\_context} \(\in \{\code{Load},\code{Store}\}\); plus operator,
boolean, unary, and comparison enumerations including bitwise forms.

\subsection{Visitor and transformer patterns}
\label{sec:visitor}

\code{NodeVisitor} dispatches on concrete types with a type-object method cache
(avoiding class-name string lookups). \code{NodeTransformer} rewrites trees
bottom-up: returning \code{None} removes a node; returning an AST replaces it;
returning a list splices into list fields~\citep{gamma1994design,asdl-python}.
Helpers \code{walk}, \code{iter\_fields}, and \code{iter\_child\_nodes} support
ad hoc analysis without subclassing. The unparser and several evaluator passes
are visitors; desugaring and future IR lowers can be transformers.

%==============================================================================
\section{AST Construction}
\label{sec:builder}
%==============================================================================

\subsection{Builder mapping: parse tree \(\to\) ASDL}
\label{sec:builder-map}

\code{PinescriptASTBuilder} extends the ANTLR-generated
\code{PinescriptParserVisitor}. Each \code{visit\_*} method maps a parser
context to one or more ASDL nodes. Design choices include:

\begin{itemize}[leftmargin=*]
  \item \textbf{Singleton operator/context nodes} (\code{Load}, \code{Store},
    \code{Add}, \ldots) to reduce allocation on hot paths.
  \item \textbf{Location attachment} from start/stop tokens, with a single-line
    fast path when the stop token contains no newline.
  \item \textbf{Store contexts} on assignment targets (names, attributes,
    subscripts, tuples).
  \item \textbf{Number and string decoding} for \code{Constant} values, including
    multiline string normalization.
  \item \textbf{Shared builder instance} across parses (stateless visits).
\end{itemize}

Deeply nested ternaries raise the recursion limit temporarily to 5000 during
walk and build.

\subsection{Annotation attachment}
\label{sec:annotations}

Annotation comments (\code{//@\ldots}) are not productions in the statement
grammar. After a successful \code{exec}-mode build, the helper:

\begin{enumerate}[leftmargin=*]
  \item Skips the entire annotation pass if the source contains no \code{@}
    (cheap filter for micro-snippets and many fixtures).
  \item Collects statement nodes via \code{StatementCollector}.
  \item Scans the token stream for \code{COMMENT} tokens containing \code{@},
    classifying kind/value via \code{PinescriptCommentParser}.
  \item Orders comments and statements by \((\mathit{lineno},\mathit{col})\),
    groups consecutive comments, and attaches them to the following node.
\end{enumerate}

Script-level annotations (kind suffix \code{S}) populate
\code{Script.annotations}; function (\code{F}), type (\code{T}), and variable
(\code{V}) annotations attach to the corresponding statement nodes.
Algorithm~\ref{alg:annot} abstracts the attachment procedure.

\begin{algorithm}[t]
\caption{AnnotationAttachment}
\label{alg:annot}
\begin{algorithmic}[1]
\Require Script AST \(T\), source text \(S\), token stream \(\tau\)
\Ensure Annotations attached on \(T\) and eligible statements
\If{\(``@'' \notin S\)}
  \State \Return \(T\) \Comment{no annotation comments possible}
\EndIf
\State \(\mathit{stmts} \gets \mathrm{StatementCollector}(T)\)
\State \(\mathit{cmts} \gets \{\,c \in \tau \mid c.\mathrm{type}=\mathrm{COMMENT} \land ``@" \in c.\mathrm{text}\,\}\)
\State \(\mathit{items} \gets \mathrm{sort}(\mathit{cmts} \cup \mathit{stmts}\ \mathrm{by}\ (\mathit{line},\mathit{col}))\)
\State Group consecutive comments vs.\ statements
\State Attach script-level \code{@*S} comments to \(T.\mathrm{annotations}\)
\For{each (comment group \(G\), statement \(s\)) pair}
  \State Attach \(G\) filtered by kind suffix to \(s.\mathrm{annotations}\) if applicable
\EndFor
\State \Return \(T\)
\end{algorithmic}
\end{algorithm}

Round-trip emits these annotation strings as full lines before the relevant
construct, preserving version directives and documentation markers used by the
LSP hover surface.

%==============================================================================
\section{Unparsing and Formatting}
\label{sec:unparse}
%==============================================================================

\subsection{Pretty-printing model}
\code{NodeUnparser} walks the ASDL tree and appends source fragments to a list
buffer. Operator precedence is tracked with a \code{Precedence} lattice so
parentheses appear only when required for associativity or binding strength.
Binary operators assign the next-lower precedence on the right-hand child to
keep left-associative chains free of redundant parentheses.

\paragraph{Normalization policy.}
The round-trip goal is \emph{semantic and structural} fidelity, not
byte-identical reproduction:

\begin{itemize}[leftmargin=*]
  \item Indentation is four spaces; operators are spaced consistently.
  \item Free-floating non-annotation comments may move or drop.
  \item Multiline string constants prefer triple-quoted forms when they contain
    newlines.
  \item Annotation lists on script/function/type/assign nodes are re-emitted.
\end{itemize}

Figure~\ref{fig:roundtrip} depicts the intended commutative diagram for
formatters and tests.

\begin{figure}[t]
\centering
\begin{tikzpicture}[
  font=\small,
  box/.style={draw, rounded corners=2pt, align=center, inner sep=6pt,
    minimum width=2.4cm, fill=codebg},
  arr/.style={-{Latex}, thick}
]
\node[box] (s1) {Source \(s\)};
\node[box, right=3.2cm of s1] (t1) {AST \(T\)};
\node[box, below=1.6cm of t1] (s2) {Source \(s'\)};
\node[box, below=1.6cm of s1] (t2) {AST \(T'\)};
\draw[arr, pyneblue] (s1) -- node[above]{\scriptsize parse} (t1);
\draw[arr, pynegreen] (t1) -- node[right]{\scriptsize unparse} (s2);
\draw[arr, pyneblue] (s2) -- node[below]{\scriptsize parse} (t2);
\draw[arr, pyneorange, dashed] (s1) -- node[left, align=center]
  {\scriptsize structural\\[-1pt]\scriptsize equivalence} (t2);
\draw[arr, pynegray, dashed] (t1) -- node[above, sloped]{\scriptsize \(\approx\)} (t2);
\end{tikzpicture}
\caption{Round-trip commutative diagram. Formatters implement
\(s \mapsto s' = \mathrm{unparse}(\mathrm{parse}(s))\); tests require
\(\mathrm{parse}(s')\) to match \(\mathrm{parse}(s)\) up to agreed normal forms
(and often check re-parse stability of \(s'\)).}
\label{fig:roundtrip}
\end{figure}

\subsection{Performance engineering}
Unparse micro-benchmarks on a 99-script corpus showed dispatch and buffer
construction dominating. Optimizations include:

\begin{itemize}[leftmargin=*]
  \item \textbf{Thread-local reuse} of a single \code{NodeUnparser} (resets
    buffers between calls; keeps type\(\to\)method cache warm).
  \item \textbf{Type-keyed dispatch} instead of class-name strings.
  \item \textbf{Precomputed indent prefixes} and operator token strings.
  \item \textbf{Lightweight context managers} for delimiters/blocks (avoiding
    \code{contextlib} generator overhead on every parenthesized subexpression).
\end{itemize}

These changes yielded approximately \(1.95\times\) higher unparse throughput on
the corpus loop (Section~\ref{sec:eval}) with byte-identical output relative to
the pre-optimization unparser.

\subsection{LSP formatting}
The language server exposes document formatting as parse \(\to\) unparse,
reusing the same path as the CLI \code{format} and \code{parse-and-unparse}
commands~\citep{lsp2023,pynescript2026}. Because annotations survive the
builder post-pass, version headers remain at the top of formatted buffers.

%==============================================================================
\section{Static Analysis}
\label{sec:static}
%==============================================================================

\subsection{Linter rules}
\label{sec:linter}

\code{PineLinter} combines a full parse with lightweight regex heuristics.
Rule codes use prefixes \(\mathrm{E}\) (errors), \(\mathrm{W}\) (warnings),
and \(\mathrm{C}\) (style). Table~\ref{tab:lint} lists the built-in rules
(nine-plus checks spanning syntax, version, deprecation, naming, and style).

\begin{table}[t]
\centering
\caption{Built-in linter rules (structural and style).}
\label{tab:lint}
\small
\begin{tabular}{@{}clp{0.55\textwidth}@{}}
\toprule
Code & Severity & Description \\
\midrule
E001 & error & Syntax error from \code{parse} (with line/column when known) \\
W001 & warning & Missing \code{//@version} declaration \\
W002 & warning & Version \(< 5\) deprecated; suggest v5/v6 \\
W101 & warning & Legacy \code{security(} pattern; prefer \code{request.security} \\
W102 & warning & Histogram plot style suggestion \\
W103 & warning & \code{var int x = na} initialization style \\
C001 & style & \code{ta.*} assignment naming heuristic (camelCase) \\
C002 & style & Line length \(> 120\) \\
C003 & style & Single-line indented \code{if} style caution \\
C004 & style & File should end with newline \\
\bottomrule
\end{tabular}
\end{table}

The linter is intentionally not a full type checker: it provides fast editor
feedback and CI gates, while deeper type modeling lives in
\file{type\_system.py}.

\subsection{Type system overview}
\label{sec:types}

The type model represents:

\begin{description}[style=nextline,leftmargin=1.2em]
  \item[Qualifiers]
    \code{TypeQualifier}: \code{const}, \code{simple}, \code{series},
    \code{input}.
  \item[Builtin scalars]
    \code{int}, \code{float}, \code{bool}, \code{string}, \code{color},
    \code{na}, and v6 \code{enum}.
  \item[Collections]
    \code{ArrayType}, \code{MatrixType}, and map-like structures with element
    (and key/value) parameters.
  \item[UDTs]
    \code{UserDefinedType} with fields and method resolution via
    \code{MethodResolver}.
  \item[Registry]
    \code{TypeRegistry} for name lookup during evaluation and tooling.
\end{description}

Inference is partial and runtime-assisted: the interpreter uses the model for
compatibility checks and UDT instantiation, while the LSP consumes related
metadata inventories (hundreds of builtins---on the order of \(482\)--\(800+\)
surface entries depending on inventory generation) for completion and
signatures~\citep{pynescript2026}. Full bidirectional checking of every Pine
qualifier lattice edge remains future work (Section~\ref{sec:limitations}).

%==============================================================================
\section{Corpus Engineering}
\label{sec:corpus}
%==============================================================================

\subsection{Sanitization of scraped sources}
\label{sec:sanitize}

Community scripts scraped from publications, forums, or documentation pages
rarely arrive as pure Pine. \file{corpus\_sanitize.py} implements
\code{sanitize\_corpus\_source}, which strips page chrome while preserving
in-comment markdown used for \code{//@function} documentation.

Typical removals include: markdown fences (\code{```pine}); blockquote chrome;
``Expand (\(N\) lines)'' UI stubs; horizontal rules; bare URLs; HTML comments;
publication footers (author/license/tags); leading prose before the first
script-like line; and mis-collected shell/Python/HTML fragments.
If no usable Pine remains, a minimal parseable stub is returned so callers fail
softly.

\begin{algorithm}[t]
\caption{SanitizeCorpus}
\label{alg:sanitize}
\begin{algorithmic}[1]
\Require Raw text \(R\) (possibly mixed markdown / UI chrome)
\Ensure Sanitized source \(S\) or minimal stub
\State \(L \gets \mathrm{splitlines}(R)\); track open quote state across lines
\State Drop fence lines, expand stubs, HR, URL-only, footer labels, HTML noise
\State Drop leading prose until a Pine-like line (\code{//@version}, 
       \code{indicator(}, \code{strategy(}, assignments, \ldots)
\State Repair mild truncation artifacts when safe
\If{\(L\) has no usable Pine}
  \State \Return minimal stub \code{//@version=5 / indicator / plot}
\EndIf
\State \Return \(\mathrm{join}(L)\)
\end{algorithmic}
\end{algorithm}

The compiler and multi-run hosts invoke sanitization on cache miss paths so warm
hits skip the cost~\citep{pynescript2026}.

\subsection{Parse-fail buckets and regeneration workflow}
\label{sec:regen}

Corpus tests parametrize over directories of \code{.pine} files. Failures are
bucketed by root cause (lexer indent, missing construct, sanitize residue,
version-specific syntax). The engineering loop is:

\begin{enumerate}[leftmargin=*]
  \item Reproduce with \code{parse}/\code{unparse} on the failing file.
  \item Prefer grammar or lexer-base fixes in \file{resource/*.g4} /
    \file{LexerBase.py}; never hand-edit \file{generated/}.
  \item Regenerate parser/lexer; align builder \code{visit\_*} names with new
    rule contexts.
  \item Extend ASDL only when node shape must change; regenerate node classes.
  \item Add a focused unit test, then re-run corpus and linter suites.
\end{enumerate}

This discipline keeps the approximate grammar evolving with community patterns
without forking generated code.

%==============================================================================
\section{Evaluation}
\label{sec:eval}
%==============================================================================

\subsection{Correctness: coverage and round-trip}
\label{sec:eval-correct}

\begin{itemize}[leftmargin=*]
  \item \textbf{Parser coverage.}
    High success rates on \(138+\) real scripts in parse-and-unparse suites
    (e.g., \(138\) passed on a builtin-script directory configuration reported
    in performance agents). Broader project automation exceeds \(1100\) tests
    across front-end, runtime, and compiler concerns~\citep{pynescript2026}.
  \item \textbf{SLL \(\equiv\) LL.}
    On \(30\) size-stratified scripts, AST dumps with attributes matched between
    SLL-first and forced pure-LL parses (\(0\) mismatches in that sample).
  \item \textbf{Annotations.}
    Scripts with \code{//@} retain \code{script.annotations} and declaration
    annotations through parse; unparse re-emits them as lines.
  \item \textbf{Modes.}
    \code{mode="eval"} parses isolated expressions; \code{mode="exec"} parses
    full scripts.
\end{itemize}

Algorithm~\ref{alg:roundtrip} states the structural check used in CI-oriented
workflows.

\begin{algorithm}[t]
\caption{RoundTripCheck}
\label{alg:roundtrip}
\begin{algorithmic}[1]
\Require Source \(s\), optional sanitize flag
\State \(s_0 \gets \mathrm{sanitize}(s)\) if enabled else \(s\)
\State \(T \gets \mathrm{parse}(s_0)\)
\State \(s' \gets \mathrm{unparse}(T)\)
\State \(T' \gets \mathrm{parse}(s')\)
\State Assert \(\mathrm{dump}(T)\) compatible with \(\mathrm{dump}(T')\)
        (structure; attributes per test policy)
\State Optionally assert selected annotation strings preserved
\end{algorithmic}
\end{algorithm}

\subsection{Parse performance}
\label{sec:eval-parse}

Profiling showed ANTLR adaptive prediction dominating wall time, with the
builder a smaller fraction. Table~\ref{tab:parse-perf} summarizes reported
measurements from the project's parse performance agents (Python~3.14,
in-process microbenchmarks).

\begin{table}[t]
\centering
\caption{Parse performance (representative agent results).}
\label{tab:parse-perf}
\small
\begin{tabular}{@{}lrrr@{}}
\toprule
Workload & Before / LL-only & After / SLL-first & Speedup \\
\midrule
Mean per script (12-file mix) & \(\sim\)80--140\,ms & \(\sim\)14.5--15\,ms & \(\sim 5.4\times\) \\
Ops/s (mix) & \(\sim\)7--13 & \(\sim\)69 & \(\sim 5\)--\(10\times\) \\
Complex \(\sim\)16\,KB script & \(\sim\)1917\,ms & \(\sim\)254\,ms & \(\sim 7.5\times\) \\
\bottomrule
\end{tabular}
\end{table}

Fair same-process A/B (12 iterations \(\times\) 12 scripts) reported approximately
\(175.7\)\,ms vs.\ \(953.7\)\,ms per pass (\(5.43\times\), \(81.6\%\) improvement)
for SLL-first vs.\ LL-only, with identical trees on success.
Additional micro-optimizations (annotation skip without \code{@}, recursion-limit
guard, location fast path, lexer deque pending tokens) further reduced constant
factors after the SLL win.

Figure~\ref{fig:sll-ll} visualizes the headline SLL vs.\ LL comparison.

\begin{figure}[t]
\centering
\begin{tikzpicture}
\begin{axis}[
  ybar,
  bar width=18pt,
  ylabel={ms / pass (12 scripts)},
  symbolic x coords={LL-only,SLL-first},
  xtick=data,
  ymin=0, ymax=1100,
  nodes near coords,
  nodes near coords align={vertical},
  width=0.72\textwidth,
  height=6.2cm,
  axis lines=left,
  enlarge x limits=0.35,
  tick label style={font=\small},
  label style={font=\small},
  every node near coord/.append style={font=\scriptsize},
]
\addplot[fill=pynegray!60] coordinates {(LL-only,953.7) (SLL-first,175.7)};
\end{axis}
\end{tikzpicture}
\caption{Parse performance: pure LL vs.\ SLL-first with LL fallback on a
12-script stratified mix (fair in-process A/B). Trees matched on the measured
sample.}
\label{fig:sll-ll}
\end{figure}

\subsection{Unparse performance}
\label{sec:eval-unparse}

On 99 scripts from an indicators/strategies set, thread-local unparser reuse and
dispatch optimizations improved median throughput from \(\sim 1266\) to
\(\sim 2469\) unparses/s (\(\approx 1.95\times\)), with large-script calls
improving by \(\approx 2.12\times\). Outputs remained byte-identical to the
baseline unparser across successfully parsed ASTs.

\subsection{Parse cache}
\label{sec:eval-cache}

A process-local LRU keyed by \((\mathrm{SHA256}(\mathit{source}), \mathit{mode})\)
(default max \(128\), \code{threading.RLock}, env toggles
\code{PYNE\_PARSE\_CACHE} / \code{PYNE\_PARSE\_CACHE\_MAX}) returns shared AST
identity on hit. On hits, evaluator-bound \code{\_pine\_call\_site} attributes
are scrubbed so multi-run hosts do not reuse another evaluator's bound methods.
Failed parses are not cached. Multi-parse warm paths report order-of-magnitude
gains approaching \(1000\times\) vs.\ cold ANTLR parse when the cache hits
repeatedly on identical sources (host benchmarks; dominated by hash + dict
lookup).

%==============================================================================
\section{Related Work}
\label{sec:related}
%==============================================================================

\paragraph{Parser generators and ANTLR.}
ANTLR~4 popularized adaptive \(\mathrm{LL}(*)\) parsing with automatic
left-recursion elimination and continuous lookahead~\citep{parr2013antlr,parr2011allstar,parr1995antlr}.
The SLL-first/LL-fallback pattern is recommended practice for throughput without
sacrificing generality; \pyne{} confirms large wins on a charting DSL with
indent-sensitive lexing layered outside pure grammar rules.

\paragraph{ASDL and language IRs.}
Zephyr ASDL~\citep{wang1997asdl} and Python's \code{ast} module~\citep{asdl-python}
demonstrate the value of declarative node schemas for compilers and tools.
\pyne{} reuses this philosophy for a DSL, including location attributes and
visitor/transformer APIs familiar to Python tooling authors.

\paragraph{DSL engineering.}
Fowler~\citep{fowler2010dsl} and Hudak~\citep{hudak1996dsl} survey external vs.\
embedded DSLs. \pine{} is an external, product-hosted language; \pyne{} builds an
external toolchain around an approximate grammar rather than embedding Pine in
Python syntax. Classic compiler texts~\citep{aho2006dragon,cooper2012engineering}
and Nystrom's interpreter construction~\citep{nystrom2021crafting} inform the
pipeline staging (lex, parse, AST, analysis, unparse).

\paragraph{Language servers and formatters.}
LSP~\citep{lsp2023} standardizes editor integration; formatters based on
parse--print pipelines (go\-fmt style) motivate lossless-enough unparsing.
\pyne{}'s format path is intentionally the same as the test round-trip.

\paragraph{Financial DSLs.}
Vendor documentation~\citep{tradingview-pine} specifies user-facing semantics but
does not publish an official grammar. Independent implementations must reverse-
engineer surface syntax from documents and corpora---a central motivation for
this paper's grammar-engineering narrative.

%==============================================================================
\section{Limitations}
\label{sec:limitations}
%==============================================================================

\begin{enumerate}[leftmargin=*]
  \item \textbf{Approximate, unofficial grammar.}
    The \file{.g4} sources are reverse-engineered and corpus-hardened. They are
    not an official TradingView grammar and may accept or reject edge cases
    differently from the proprietary compiler.
  \item \textbf{Semantic version gaps.}
    Some v6 semantic rules (strict booleans, division typing) are enforced in
    later stages, not fully in the parser.
  \item \textbf{Comment fidelity.}
    Only \code{//@}-style annotations are first-class on the AST; arbitrary
    comments are not a full comment-preserving CST.
  \item \textbf{Linter depth.}
    Rules are heuristic; the type model is not a complete static type checker
    for the Pine qualifier lattice.
  \item \textbf{Cache mutability.}
    Shared AST identity requires read-only treatment or explicit
    \code{clear\_parse\_cache} after transformations.
  \item \textbf{Corpus bias.}
    Evaluation scripts under-represent some rare library and brokerage-specific
    patterns; sanitize stubs can mask total scrape failure.
\end{enumerate}

%==============================================================================
\section{Conclusion}
\label{sec:conclusion}
%==============================================================================

We described the \pyne{} language front-end for \pine{}: an \antlr{} grammar with
Python indent tracking, a Zephyr-\asdl{} IR, a visitor builder with annotation
attachment, a precedence-aware unparser, a structural linter, and a qualifier-
aware type model, complemented by corpus sanitization and regeneration
discipline.
The pipeline enables formatters, LSP features, interpreters, and compilers to
share one tree shape. Measured optimizations---SLL-first parsing
(\(\approx 5.4\times\)), unparser reuse (\(\approx 1.95\times\)), and SHA-256
AST caching (orders of magnitude on multi-parse hits)---show that careful
front-end engineering pays off even when ANTLR already dominates profiles.

Future work includes tighter CST/comment preservation, richer static checking of
qualifiers and UDTs, grammar differential testing against larger public corpora,
and continued v6+ surface tracking as the language evolves.
\pyne{} remains an independent open implementation; we hope the documented
grammar-engineering practices help other charting-DSL and financial-DSL tooling
efforts~\citep{pynescript2026}.

\paragraph{Acknowledgments.}
We thank contributors to the \pyne{} / HOOX open-source effort and authors of
ANTLR, ASDL, and the Python \code{ast} ecosystem.

\paragraph{Trademark notice.}
\trademarknote{}

%==============================================================================
\bibliographystyle{plainnat}
\bibliography{../common/bibliography}

\end{document}
